\documentclass[11pt]{article}
\usepackage{amsmath,amssymb,amsthm,hyperref}
\usepackage[margin=1in]{geometry}
\newtheorem{lemma}{Lemma}
\newtheorem{theorem}{Theorem}
\newtheorem{corollary}{Corollary}
\begin{document}
\title{Ordinary support $\{3,7\}$ by a Frobenius minor}
\author{}\date{}\maketitle
For each integer $n\ge1$, we prove consecutive coprimality for
$P_n(t)=\sum_{k=0}^n\binom nk(3n+1)_{4k}t^k$.
The earlier residue reduction leaves three progressions.
A restricted power-series minor controls every common-root candidate
on each remaining index. Rational conjugacy and even-moment
positivity exclude the last possible root.

The exact inputs are Lemma \texttt{lem:mixed-matrix},
Corollary \texttt{cor:sixty-one-sixty-four}, and the strict-radius
argument in the proof of Corollary \texttt{cor:fifteen-sixteenths}
of the \href{LOWER_THREE_FOURTH_POWER_RESIDUE_REDUCTION.tex}{residue note};
and Lemma \texttt{lem:artin-hasse-closure} of the
\href{ARITHMETIC_FRONTIER_TERMINAL_ORDERS.tex}{arithmetic note}.
These are applied only to the actual fixed-shape auxiliaries,
at ordinary integer indices.

Work in a complete algebraically closed extension of $\mathbb Q_2$,
with $v=v_2$, $v(2)=1$, valuation ring $\mathcal O$, maximal ideal
$\mathfrak m$ and residue field $k=\mathcal O/\mathfrak m$. Congruences of restricted series below are
coefficientwise over $\mathbb Z_2$.
For the minor construction fix $n\ge1$ in one of the three remaining classes.
Write
\[
 \gamma=3n+1,\ c=3n+4,\ M=7n+1,\ m=M+11,\quad
 K=7(7n+5)(7n+6),\quad E=v_2(K).
\]
Put $R_i(t)=\prod_{\ell=n}^{n+i-1}(t+4\ell)$ and
$\Pi_{n,q}(t)=\prod_{j=0}^{n-1}(t+qj)$.
All auxiliaries use the same $\gamma$, with raising operator
$A=c-zD_z+z$ and $B_j=\Pi_{j,4}(A)1$.
We abbreviate $B=B_n$, $C=B_{n+1}$ and $D=B_{n+2}$.
The imported exact matrix gives
\[
 \operatorname{rank}[B_{n+1},B_{n+2}-4uB_n/K,
              B_{n+3}-4muB_n/K](u)<3                 \tag{1}
\]
at a common root, where $u=-1/t$.
On every remaining progression every common root has $v_2(u)>3$.
For odd remaining $n$, $E\ge5$; for even remaining $n$,
$n\equiv54\pmod{64}$ and $E\ge6$.

\begin{lemma}[A Frobenius minor for every remaining index]
\label{lem:frobenius-minor}
Fix $n\equiv13,45,54\pmod{64}$.
There is a restricted series $F_n\in\mathbb Z_2\langle w\rangle$
such that every actual common root of $P_n,P_{n+1}$ gives
$w=-1/(8t)$ in the maximal ideal and $F_n(w)=0$.
For odd $n$, $\bar F_n=w$.
For even $n$, $F_n$ has unit linear coefficient.
\end{lemma}
\begin{proof}
\par\medskip\noindent\textit{A fixed polynomial column basis.}
Let $w=u/8$ and $a=K/32$, so $a\in\mathbb Z_2$ and $v(w)>0$.
Invertible scalar column operations reduce (1) to the cyclic
columns of $\Pi_{n,4}(A)T_i(A)1$, where
\[
 T_1(t)=t+4n,\quad T_2(t)=aR_2(t)-w,\quad
 T_3(t)=(t+4n)(t-1)(t+n+1).                         \tag{2}
\]
To check this, clear the second original column to $KD-4uB$.
Replace the third by $B_{n+3}-mB_{n+2}$, whose polynomial is
$R_2(t)(t-c)$, and divide the second by32.
Finally subtract $\lambda T_1$ from the third, with
$\lambda=(4n+5)(1-c)$.
The factor identity is
\[
 R_2(t)(t-c)-\lambda(t+4n)
             =(t+4n)(t-1)(t+n+1).
\]
The outer parameter $w$ is an evaluated scalar.
All Euler operators act on the inner variable only.

\par\medskip\noindent\textit{The exact two-block interface.}
Apply the Artin--Hasse stage change at $p=2,q=4$, with original
seed1 and no higher multiplication terms.
Its inner variable is $x=-z^2/2$, so $x^2=u/4=2w$.
The nonresonant block has operator $1+2L_N$; its integral seed,
after absorption of $\Pi_{n,4}(1+2L_N)$, is denoted $h_N$.
The resonant block has operator $2L_R$, and seed $h_R$:
\[
 L_j=c_j-xD_x+x-2x^2-8x^4-\cdots,\qquad c_j\in\mathbb Z_2.
\]
More explicitly, if $E_2(z)^{-1}=\sum e_jz^j$, put
$H_j(x)=\sum_{\ell\ge0}e_{2\ell+j}(-2x)^\ell$.
Choose $j_R\in\{0,1\}$ congruent to $c$ modulo2 and $j_N=1-j_R$.
Then
\[
 c_R=(c-j_R)/2,\qquad c_N=(c-j_N-1)/2,\qquad
 h_R=H_{j_R},\quad
 h_N=\Pi_{n,4}(1+2L_N)H_{j_N}.
\]
Both constants are integers and both seeds belong to
$\mathbb Z_2[[x]]$ with constant reduction1.
For odd $n$, $c_R=c_N=(3n+3)/2$.
For even $n$ in the remaining class,
\[
 c_R=c/2=(3n+4)/2\equiv3\pmod4,\qquad c_N=c_R-1\equiv2\pmod4.
                                                               \tag{3}
\]
The exact block columns are
\[
 {\cal C}_2(T_i(1+2L_N)h_N;2w),\qquad
 2^n{\cal C}_2(\Pi_{n,2}(L_R)T_i(2L_R)h_R;2w).       \tag{4}
\]
A zero dependence of the original columns gives the same zero
dependence in both blocks.

Use the ring $\mathbb Z_2\langle w\rangle$ of restricted power
series, whose coefficients are integral and tend to zero.
For an integral inner series independent of the outer parameter,
its cyclic coefficient at outer degree $\ell$ has a factor $2^\ell$
from $x^2=2w$. Finite Euler operations and integral inner seeds
preserve inner integrality, so these vectors belong to this ring.
Multiplication by the integral outer polynomial coefficients of
$T_i$ preserves restricted-series integrality and convergence. Thus no characteristic-zero Euler action on a
fixed nonzero cyclic quotient is being assumed.

In the nonresonant block set
\[
 v_0={\cal C}_2(h_N;2w),\qquad
 v_1={\cal C}_2(L_Nh_N;2w).
\]
These form a unit basis over $\mathbb Z_2\langle w\rangle$.
Indeed their reduction is the constant cyclic basis for
$x-xD_x+\bar c_N$ on $k[x]/(x^2)$.
The determinant has odd constant reduction and every other reduced
coefficient zero. Its inverse is integral restricted, by the
geometric series for a unit of the form $b+2H(w)$.
Define coordinates
\[
 {\cal C}_2(L_N^2h_N;2w)=\alpha v_0+\beta v_1,\quad
 {\cal C}_2(L_N^3h_N;2w)=\alpha_3v_0+\beta_3v_1.
\]
The idempotent reduced operator gives
\[
 \alpha,\alpha_3\in2\mathbb Z_2\langle w\rangle,\qquad
 \bar\beta=\bar\beta_3=1.                              \tag{5}
\]
After conversion to this basis the second nonresonant coordinate
row is divisible by2 for each column in (2), since
$T(1+2L)$ has divided-Taylor coefficients $2^rT^{[r]}(1)$.
Divide that row by2; call the resulting two rows $N_0,N_1$.
With $r=4n+1$, their entries are exactly
\[
 \begin{array}{c|ccc}
 &T_1&T_2&T_3\\ \hline
 N_0&r&a r(r+4)-w+4a\alpha&
             4(5n+3)\alpha+8\alpha_3\\
 N_1&1&2a(r+2+\beta)&
             r(n+2)+2(5n+3)\beta+4\beta_3 .
 \end{array}                                           \tag{6}
\]
All these divisions are by integer powers of2 and all resulting
entries are integral restricted series.

\par\medskip\noindent\textit{Odd indices.}
Suppose $n$ is odd. The reduction of (6) is
\[
 \begin{bmatrix}1&\bar a-w&0\\1&0&1\end{bmatrix}.
\]
Remove $2^n$ from the resonant rows.
The first and third columns are even there, because
\[
 T_1(2L)=2(L+2n),\quad
 T_3(2L)=2(L+2n)(2L-1)(2L+n+1),
\]
and $n+1$ is even.
The second column reduces to $-w$ times
$\Omega={\cal C}_2(\Pi_{n,2}(L_R)h_R;2w)$.
The vector $\bar\Omega=\bar L_R^n1$ is nonzero,
as $\bar L_R^2=\bar L_R$ and $\bar L_R1\ne0$.
Choose a component with reduction1.
The determinant $G(w)$ of these three rows is integral restricted,
and
\[
 \bar G(w)=
 \det\begin{bmatrix}1&\bar a-w&0\\1&0&1\\0&w&0\end{bmatrix}=w.
                                                               \tag{7}
\]
Consequently $G$ is injective on the integral disk: its linear
coefficient is a unit and all higher coefficients are even.

\par\medskip\noindent\textit{The even remaining progression.}
Now $n\equiv54\pmod{64}$.
In particular $n\equiv6\pmod8$, $a$ is even,
$r\equiv1\pmod4$, and $d=(n+2)/2\equiv0\pmod4$.
The whole third column of the normalized four-row block matrix is
even, so divide it by2. In the nonresonant block this follows also from
\[
 T_3(1+2L)=4L(L+d)(r+2L);
\]
in the resonant block it follows from the factor2 of $T_3(2L)$.

After this column division, (5)--(6) give the first two rows modulo4:
\[
 [1,\ a-w,\ 0],\qquad [1,\ 0,\ \tau(w)],
\]
where
\[
 \tau=2\beta_3+(r+2d)\beta+rd,\qquad \bar\tau=1.       \tag{8}
\]
For example the new $N_0$ third entry is
$2[2\alpha_3+(r+2d)\alpha]$, divisible by4.
The $N_1$ second entry is zero modulo4 because $a$ is even.

Take the constant-component row of the resonant block, after
removal of $2^n$.
Write
\[
 \omega(w)={\cal C}_2(\Pi_{n,2}(L_R)h_R;2w)_0,\quad
 \nu(w)={\cal C}_2(\Pi_{n,2}(L_R)L_Rh_R;2w)_0.
\]
Both reduce to the constant1, by (3).
The three entries of this row modulo4 are
\[
 [\,2\nu,\ -w\omega,\ \nu\,].                         \tag{9}
\]
For the last entry use $2n\equiv0\pmod4$, $n+1\equiv3\pmod4$:
\[
 (L+2n)(2L-1)(2L+n+1)\equiv L\pmod4.
\]

Let $G(w)$ be the determinant of these three normalized rows.
Equations (8)--(9), as identities of restricted-series coefficients,
give
\[
 G(w)\equiv a\nu(w)+w\bigl(\tau(w)\omega(w)-\nu(w)\bigr)\pmod4.
                                                               \tag{10}
\]
Its reduction is zero, so $H=G/2$ is integral restricted.
Also $\bar\nu=1$ and $\tau\omega-\nu\in2\mathbb Z_2\langle w\rangle$.

We compute the linear coefficient of $H$ modulo2 at $w=0$.
The coefficient of $w$ in $\nu(w)$ is even, since it arises from
the substitution $x^2=2w$. At $w=0$, the first-two-coefficient
operator for $L_N$ has diagonal $c_N,c_N-1$; it preserves
the quotient by $x^2$. Cayley--Hamilton therefore gives
\[
 \beta(0)=2c_N-1,\qquad \beta_3(0)\equiv1\pmod2.
\]
Using (3) and (8), $\tau(0)\equiv1\pmod4$.
The constant coefficient of $L_Rf$ is $c_R f(0)$, so
\[
 \nu(0)=c_R\omega(0),\qquad \omega(0)\in\mathbb Z_2^\times.
\]
Equation (10) now shows
\[
 [w]H\bmod2
   =\frac{\tau(0)\omega(0)-\nu(0)}2\bmod2
   =\frac{1-c_R}{2}\omega(0)\bmod2=1.                \tag{11}
\]
Thus $H$ has unit linear coefficient and integral higher coefficients.
It is injective on the maximal ideal: at two distinct arguments there,
every higher term in its divided difference has positive valuation.
Convergence follows from the restricted-series condition.

Take $F_n=G$ for odd $n$, and $F_n=G/2$ for even $n$.
\end{proof}

\begin{theorem}[Every ordinary consecutive pair]\label{thm:coprime}
For every integer $n\ge1$, $\gcd(P_n,P_{n+1})=1$ over
$\mathbb Q$ and over $\mathbb C$.
\end{theorem}
\begin{proof}
The retained residue theorem covers all other indices.
On each remaining progression use Lemma~\ref{lem:frobenius-minor}.
For odd $n$, the series with reduction $w$ is injective on the
integral disk: its divided difference has unit constant term and
all other terms even.
For even $n$, its linear coefficient is a unit and its higher
coefficients are integral. At two distinct arguments in the
maximal ideal every higher term in its divided difference has
positive valuation. The series converges, so this difference
quotient is again a unit.
Thus there is at most one distinct common root over $\mathbb C_2$.

All conjugates over $\mathbb Q$ of an actual rational common factor
obey the same rank identities and disk bound. Characteristic-zero
separability would force that factor to have degree one.
Its root would be real. But for even $j$ and real $t$,
\[
 P_j(t)=\frac1{\Gamma(3j+1)}
 \int_0^\infty x^{3j}(1+tx^4)^j e^{-x}\,dx>0.
\]
The integrand is nonnegative and positive on an interval.
One of the two indices is even, a contradiction.
A rational Bezout identity extends coprimality to $\mathbb C$.
This uses conjugates of an actual common factor, not rationality
of an arbitrary root of an auxiliary restricted series.
\end{proof}

\begin{corollary}\label{cor:sfc}
Every nonzero complex phase $X^3(A+BX^4)$ satisfies ordinary
Strong Factorial at every consecutive positive pair.
\end{corollary}
\begin{proof}
For $AB\ne0$ its order-$j$ moment is $A^j(3j)!P_j(B/A)$,
so the theorem applies. If a coefficient vanishes, the nonzero
phase is a monomial and every positive factorial moment is nonzero.
\end{proof}

\paragraph{Evidence and scope.}
This proves the entire ordinary support $\{3,7\}$.
Arbitrary positive Gamma shapes, other degree-seven supports and
unrestricted binomial or arbitrary-support SFC are separate claims.
The checker reconstructs actual finite Artin--Hasse block jets,
verifies symbolic column/Taylor identities and coefficientwise
modulo-four calculations, and checks normalization integrality.
The restricted-series proof supplies the unbounded quantifiers.
Internal audit and replay are distinct from independent
whole-proof replay, formal verification, external review and priority.
\end{document}
